\section{Lowering operators and hook-entropy moments}
\label{ten:lowering}
We prove two estimates for the sweep $B_n$, $n=2^d$. The first uses
ordinary tensor lowerings to bound the operator norm in terms of the
first-row defect $k=n-\lambda_1$. The second bounds an unnormalized
moment using the best hook subdiagram of $\lambda$. It uses the first
estimate at sparse levels, then a separate signed-tensor interpolation
at the remaining levels.

\subsection{The two quantitative bounds}

\begin{proposition}[A norm bound from lowering]\label{s:lowering-norm}
There are absolute constants $c>0$ and $C$ such that, for every
dyadic $n\ge2$ and every $\lambda\vdash n$ with
$2\le k=n-\lambda_1\le n/2$,
\begin{equation}
 \|B_n\|_\lambda^2 \le \exp(-c k\log(n/k)+C k),\qquad
                         2\le k\le n/2                     \label{ten:lowering:eq:1}
\end{equation}
The block with defect one is zero.
\end{proposition}

The bound is valid throughout the displayed range; it gives a
contraction when $c\log(n/k)>C$.

Here is the second target. Fix
\[
 \beta=\tfrac1{10},\qquad \delta=10^{-3},\qquad \alpha=10^{-6},
\]
with $\alpha$ decreased further if needed, and set
$q_n=\lfloor n^\delta\rfloor$. A $q$-hook diagram has no cell
$(q+1,q+1)$. For a diagram $\eta$ of size $N$, let $a$ be its
Durfee index and use its Frobenius parts
\[
 a_i=\eta_i-i+1,\qquad b_i=\eta'_i-i\qquad(1\le i\le a).
\]
These parts sum to $N$. Define
\[
 F(\eta)=\sum_{i:a_i>0}a_i\log(N/a_i)
          +\sum_{i:b_i>0}b_i\log(N/b_i),
 \qquad F(\varnothing)=0,
\]
and the clipped remainder cost
\begin{equation}\label{ten:lowering:eq:3}
 G_n(t)=t\max\{0,\alpha\log n+\log(t/n)\}\quad(0<t\le n),
 \qquad G_n(0)=0.
\end{equation}

\begin{proposition}[Hook-entropy moment bound]
\label{ten:lowering:main}
For the fixed constants above, there is a family of real powers
$p_n\ge8$, bounded uniformly over dyadic $n\ge2$, such that every
$\lambda\vdash n$ satisfies
\begin{equation}\label{ten:lowering:eq:6}
 D_\lambda\operatorname{Tr}_\lambda |B_n|^{p_n}
 \le \exp L_n(\lambda),\qquad
 L_n(\lambda)=\min_{\substack{\eta\subseteq\lambda\\q_n\text{-hook}}}
             \{\beta F(\eta)+G_n(n-|\eta|)\}.
\end{equation}
\end{proposition}

The minimum includes the empty diagram. Equivalently, the moment
is bounded by the cost of every permitted hook subdiagram. This
quantifier will allow us to use the inherited subdiagram in each
child block, after shrinking its hook parameter when necessary.
The moment proof uses the norm estimate at
$k\le n^{1-\alpha/16}$. All notation in the rest of this section
is local; traces and tensor norms remain unnormalized.

\subsection{Ordinary tensors and the backward split}

We first prove Proposition~\ref{s:lowering-norm}. Split the sites
into two halves of size $m=n/2$. The intermediate object will be a
probability law on the half excitation counts $l_i$ and half defects
$j_i$, $i=0,1$. We will show that $l_0$ is concentrated around $k/2$
and that the loss $k-j_0-j_1$ has a small upper tail. Those two
facts will close the norm induction.

The tensors in this part carry the ordinary permutation action.
The signed color action used for Proposition~\ref{ten:lowering:main}
will be introduced after the norm proof.
All tensor words are orthonormal, so the norms below are counting norms,
with no probability normalization depending on the excitation sector.
For the sweep application fix \(n=2^d\) and
\(\lambda=(n-k,\xi)\) with \(2\le k\le n/2\).

Let \(q\ge1\) be an integer, let \(V=\mathbb C^q\), with orthonormal
basis \(e_1,\ldots,e_q\), and put
\(\mathcal A=\mathbb C e_0\oplus V\).  For a finite set \(X\) of sites write
\(\mathcal H_X=\mathcal A^{\otimes X}\), and let
\(\mathcal H_{X,l}\) be the span of words with exactly \(l\) entries in
\(V\).  A permutation \(g\) of \(X\) sends the letter at \(x\) to \(g(x)\).
In particular it introduces no signs.  We use the zero space when \(l<0\)
or \(l>|X|\).

For \(c\in\{1,\ldots,q\}\), let
\[
 a_{x,c}=|e_0\rangle\langle e_c|_x,\qquad
 L_{Y,c}=\sum_{x\in Y}a_{x,c}\quad(Y\subseteq X).
\]
Thus \(L_{X,c}:\mathcal H_{X,l}\longrightarrow\mathcal H_{X,l-1}\).
Its adjoint for the stated norms is the restriction of
\(\sum_{x\in X}|e_c\rangle\langle e_0|_x\) in the reverse direction.
More generally its half version has source and target
\[
 L_{Y,c}:\mathcal H_{Y,l}\otimes\mathcal H_{X\setminus Y,r}
 \longrightarrow
 \mathcal H_{Y,l-1}\otimes\mathcal H_{X\setminus Y,r},
\]
with adjoint \(\sum_{x\in Y}|e_c\rangle\langle e_0|_x\) in the
reverse direction.  Conjugation gives
\(\rho(g)L_{Y,c}\rho(g)^{-1}=L_{gY,c}\).  All the \(a_{x,c}\)
commute: on different
sites they act on separate factors, while on one site either product of
two of them is zero.  Hence all lowerings just defined commute.  The global case is the
standard diagonal matrix-unit action~\cite[Theorem~4.55]{etingof}.

The norms of these maps are also explicit.  On a nonzero displayed
sector, so $0\le r\le|X\setminus Y|$, and for $1\le l\le |Y|$,
\[
 \|L_{Y,c}\|_{\mathrm{op}}^2=l(|Y|-l+1),
\]
and the map is zero at $l=0$.  To see this, fix the sites and values
of all letters other than $0,c$.  If $h$ sites carry $c$, the remaining
map is the incidence matrix from $h$-subsets to $(h-1)$-subsets.  It
has $h$ ones in each column and $|Y|-l+1$ in each row.  Cauchy--Schwarz
gives squared norm $h(|Y|-l+1)$, attained on constant vectors;
maximizing at $h=l$ gives the formula.  It applies to the global map
with $Y=X$ and, when $|X|=2m$, to a half map with $|Y|=m$.

Here is the relation to the tuple realization and the multiplicities that
will be retained below.  If \(|X|=N\), fixing the \(l\) excited sites gives
the ordinary induced representation
\[
 \mathcal H_{X,l}\simeq
 \operatorname{Ind}_{S_{N-l}\times S_l}^{S_N}
       \bigl(\mathbf1\otimes V^{\otimes l}\bigr).
\]
The direct sum over the excited subsets has its orthogonal tensor norm.
Ordinary Schur--Weyl duality~\cite[Theorem~4.57 and Corollary~4.59]{etingof}
and the Pieri rule, the one-row Littlewood--Richardson case~\cite[\S2.10, equation~(2.10.1)]{SamSnowden2012}, give the unitary decomposition
\begin{equation}\label{s:lowering-full-multiplicities}
 \mathcal H_{X,l}\simeq
 \bigoplus_{\mu\vdash N}V_\mu\otimes\mathcal M_{\mu,l},
 \qquad
 \mathcal M_{\mu,l}=
 \bigoplus_{\substack{\tau\vdash l,\ \ell(\tau)\le q\\
                  \mu/\tau\text{ a horizontal strip of size }N-l}}
                  \mathbf S_\tau(V).
\end{equation}
Each pair \((\mu,\tau)\) in this sum has Pieri multiplicity one.  The full
Schur color module \(\mathbf S_\tau(V)\) is retained.  The hook-content
formula~\cite[\S3.6, equation~(3.6.1), and \S4.3]{SamSnowden2012} gives
\[
 \dim\mathbf S_\tau(\mathbb C^q)
   =\prod_{(r,c)\in\tau}\frac{q+c-r}{h_\tau(r,c)}.
\]
Here \(h_\tau(r,c)\) is the hook length of the indicated cell.
At \(l=0\) this is trivial induction and the empty partition contributes
one copy of the trivial representation; at \(l=N\) it is identity
induction.  If \(\mu=(N-j,\xi)\) occurs, the horizontal-strip
interlacing inequalities give \(\tau_r\ge\mu_{r+1}=\xi_r\).
Consequently \(j=|\xi|\le l\).

For a global defect \(k\), take \(q=k\).  The subspace of
\(\mathcal H_{X,k}\) in which each of the letters \(1,\ldots,k\) occurs
once is isometric to the counting-norm space on ordered injections:
\[
 \mathbf e_{(x_1,\ldots,x_k)}\longmapsto
 \text{the word with letter \(c\) at \(x_c\) and \(0\) elsewhere}.
\]
This isometry intertwines the position actions.  More generally, for
a set $A$ of distinct excited names let $\mathcal I_{X,A}$ be the sector
in which every name in $A$ occurs once and no other excited name occurs.
For $c\in A$, the restriction
$L_{X,c}:\mathcal I_{X,A}\to\mathcal I_{X,A\setminus\{c\}}$
and its adjoint in the reverse direction act on counting-norm injection
coefficients exactly as
\[
 (L_{X,c}f)(y)=\sum_{z\in X\setminus y(A\setminus\{c\})}
    f(y\cup\{c\mapsto z\}),\qquad
 (L_{X,c}^*g)(x)=g(x|_{A\setminus\{c\}}).
\]
Thus the restricted squared norm is $N-|A|+1$, the size of each
deletion fiber; if $c\notin A$, the restriction is zero.  For a fixed
assignment of the present names to two halves of size $m$, the same
formulas hold inside a half, with squared norm $m-l+1$ when that half contains $l$
names including $c$.  This identifies the operation with raw tuple
deletion, including its adjoint and its normalization.

In \(\mathcal I_{X,[k]}\), the shape
\(\lambda=(N-k,\xi)\) occurs with the minimum-slot multiplicity
\(f^\xi\) from Section~\ref{sec:prelim}.  It occurs on no sector
\(\mathcal H_{X,l}\) with \(l<k\), also directly by
\eqref{s:lowering-full-multiplicities}.  Each \(L_{X,c}\) intertwines
\(S_N\), so it is zero on the entire \(\lambda\)-isotypic subspace of
\(\mathcal H_{X,k}\).  This proves the needed global cancellation with
the tensor adjoints and norms fixed above.

\paragraph{The backward split.}
Take \(X=\mathbb F_2^d\), write \(n=|X|=2m\), and use the last coordinate
of the fixed chronological order \(1,\ldots,d\).  Its two halves are
\(X_0,X_1\), and its matching
pairs are \(\{(p,0),(p,1)\}\), \(p=1,\ldots,m\).  Let \(P=\Pi_d\) be
the average of their independent switches.  The preceding layers act
separately in the two halves, and in their original order they give
\[
 B_n=P C,\qquad
 C=B_{m,0}\otimes B_{m,1},\qquad
 B_nB_n^*=P(CC^*)P.
\]
The tensor product is the action of the subgroup
\(S(X_0)\times S(X_1)\); inside each factor \(B_{m,i}\) has chronological
order \(1,\ldots,d-1\).  For the backward calculation the child operator
is therefore \(C^*=B_{m,0}^*\otimes B_{m,1}^*\).

Fix one of the selected \(V_\lambda\) copies in the distinct-letter
subspace just identified.
If \(\|B_n\|_\lambda=0\) there is nothing to prove.  Otherwise choose a
unit top eigenvector \(v\) of \(B_nB_n^*\) in this copy.  Its eigenvalue
is \(\|B_n\|_\lambda^2>0\), and the last displayed identity implies
\[
 Pv=v,\qquad
 \|B_n\|_\lambda^2=\|C^*v\|^2.
\]
This uses the positive square with the last layer on both outside ends;
it preserves the fixed sweep orientation.

To define the component law, let \(E_{i,l}\) be the orthogonal projection
onto words with \(l\) excitations in \(X_i\).  For \(\mu\vdash m\) let
\[
 e_\mu^{(i)}
   =\frac{D_\mu}{m!}\sum_{g\in S(X_i)}
        \chi_\mu(g^{-1})\,\rho(g)
\]
be the central orthogonal isotypic projection.  The four factors in
\[
 Q_{l_0,l_1;\mu_0,\mu_1}
   =E_{0,l_0}E_{1,l_1}e_{\mu_0}^{(0)}e_{\mu_1}^{(1)}
       \bigm|_{\mathcal H_{X,k}},
 \qquad l_0+l_1=k,
\]
commute.  These are mutually orthogonal projections summing to the
identity.  Their ranges are, up to the unitary decomposition,
\(V_{\mu_0}\otimes V_{\mu_1}\otimes
\mathcal M_{\mu_0,l_0}\otimes\mathcal M_{\mu_1,l_1}\), with the entire
multiplicity spaces from
\eqref{s:lowering-full-multiplicities}.  In particular we do not
choose one copy, divide by its dimension, or assume any distribution on
the multiplicity spaces.  Define
\[
 \Pr(l_0,l_1,\mu_0,\mu_1)
       =\|Q_{l_0,l_1;\mu_0,\mu_1}v\|^2,\qquad
 j_i=m-(\mu_i)_1,\qquad h_i=l_i-j_i .
\]
The probabilities sum to one, and \(0\le j_i\le l_i\).  Since \(C^*\)
preserves every projected range and acts there as
\(B_m^*|_{V_{\mu_0}}\otimes B_m^*|_{V_{\mu_1}}\) tensored with the
identity on the full multiplicity space,
\begin{equation}\label{s:lowering-backward-recursion}
 \|B_n\|_\lambda^2
 \le\mathbb E\!\left[
       \|B_m\|_{\mu_0}^2\|B_m\|_{\mu_1}^2\right].
\end{equation}
Thus this is a spectral weighting from \(v\), with no independence
assumption on the two shapes.

To use this recurrence, we need two properties of its weights:
\(l_0\) is concentrated around \(k/2\), and the loss
\(k-j_0-j_1=h_0+h_1\) has a small upper tail. The matching symmetry
proves the first, and the lowering comparison below proves the second.

The marginal law of \(l_0\) is the ordinary word-coordinate law:
summing the isotypic projections gives
\(\Pr(l_0=a)=\|E_{0,a}v\|^2\).  Sample a word \(w\) with probability
\(|\langle w,v\rangle|^2\).  The equality \(Pv=v\) says that \(v\) is
fixed by each pair swap, so this coordinate law is uniform on each
orbit of those swaps.  If \(t\) pairs of \(w\) have two excitations,
then \(k-2t\) pairs have one.  The doubled pairs put one excitation in
each half, and the singly occupied pairs have independent fair
orientations on that orbit.  Hence \(l_0\) is a mixture of
\[
 t+\operatorname{Bin}(k-2t,1/2).
\]
Conditionally on the orbit, for every real $z$ the centered exponential
moment is
\[
 \mathbb E\bigl[e^{z(l_0-k/2)}\mid\text{orbit}\bigr]
 =\cosh(z/2)^{k-2t}\le e^{kz^2/8}.
\]
Here $\log\cosh r\le r^2/2$ follows by integrating
$|\tanh r|\le|r|$.  Optimizing Markov's inequality for either tail gives
\begin{equation}\label{s:lowering-split-law}
 \Pr(|l_0-k/2|\ge y)\le2e^{-2y^2/k}\quad(y\ge0),
 \qquad
 \Pr(l_0=0\text{ or }l_0=k)\le2^{1-k}.
\end{equation}

\paragraph{A positive lowering square.}
We next control \(h_i\), the difference between the half excitation
count and its first-row defect.  For this local calculation let \(X\)
be a set of \(m\) sites and put \(L_c=L_{X,c}\).  Let \(P_x^0,P_x^+\) project onto the base and excited letters at \(x\).
Write \(S_{xy}=\rho((xy))\), and define
\[
 \Omega_{\mathrm{pos}}=\sum_{x<y}S_{xy},
 \qquad
 \Omega_{\mathrm{col}}=\sum_{x<y}P_x^+P_y^+S_{xy}.
\]
For the adjoints already fixed, direct multiplication on
\(\mathcal H_{X,l}\) gives the exact identity
\begin{equation}\label{s:lowering-square}
 \sum_{c=1}^q L_c^*L_c
 =lI+\Omega_{\mathrm{pos}}
       -\binom{m-l}{2}I-\Omega_{\mathrm{col}}.
\end{equation}
Indeed the same-site terms sum to \(\sum_xP_x^+=lI\).
For \(x<y\) the two cross terms
\(\sum_c(a_{x,c}^*a_{y,c}+a_{y,c}^*a_{x,c})\) exchange a zero and
an excited letter and vanish on the other two excitation sectors of
the pair.  The full site exchange is the sum of this mixed exchange,
the zero--zero exchange, and the excited--excited exchange.  The sum
of the zero--zero exchanges is \(\binom{m-l}{2}I\), which proves
\eqref{s:lowering-square}.

Let \(c(\rho)=\sum_{(r,b)\in\rho}(b-r)\).  The transposition content
formula~\cite[Eq.~(2.1), Proposition~5.3, and Theorem~5.8]{VershikOkounkov2005} makes
\(\Omega_{\mathrm{pos}}\) act by \(c(\mu)\) on \(V_\mu\).
On a fixed excited set \(T\), the restriction of
\(\Omega_{\mathrm{col}}\) to \(V^{\otimes T}\) is precisely the
ordinary transposition sum on those \(l\) tensor factors.  It therefore
acts by \(c(\tau)\) on the summand indexed by \(\tau\) in
\eqref{s:lowering-full-multiplicities}; this remains true on its
induction to all excited sets.  Thus the exact eigenvalue of
\eqref{s:lowering-square} on the \((\mu,\tau)\) summand is
\[
 l+c(\mu)-\binom{m-l}{2}-c(\tau).
\]
This statement includes the whole color space \(\mathbf S_\tau(V)\).

Write \(\mu=(m-j,\xi)\) and \(h=l-j\).  Pieri gives
\(\tau\supseteq\xi\), with \(h\) extra boxes.  The first row and shifted
tail of \(\mu\) give
\[
 c(\mu)=\binom{m-j}{2}+c(\xi)-j.
\]
Add the \(h\) boxes of \(\tau/\xi\) in any Young-diagram order.  A
box added to a diagram of size \(a\) has content at most \(a\);
therefore
\[
 c(\tau)\le c(\xi)+hj+\binom h2 .
\]
The scalar lower bound simplifies exactly as
\[
 l+\binom{m-j}{2}-j-\binom{m-l}{2}-hj-\binom h2
       =h(m-l-j+1).
\]
Thus we obtain the uniform operator inequality
on the whole \(\mu\)-isotypic subspace:
\begin{equation}\label{s:lowering-one-step}
 \sum_cL_c^*L_c
 \succeq(l-j)(m-l-j+1)I
 \succeq(l-j)(m-2l+1)I .
\end{equation}
The first inequality is sharper; the second gives the coarser form
\((m-O(l))(l-j)\).

For the rest of the lowering comparison assume \(k/m\le1/4\).
For an integer \(s\ge1\), define the ordered-letter sum in half \(i\) by
\[
 U_{i,s}(v)=
 \sum_{c_1,\ldots,c_s=1}^q
 \|L_{X_i,c_s}\cdots L_{X_i,c_1}v\|^2.
\]
Each \(L_{X_i,c}\) intertwines both half position groups.  After \(s\)
lowerings it sends the count pair \((l_i,l_{1-i})\) to
\((l_i-s,l_{1-i})\).  Hence for each fixed ordered letter list,
distinct count pairs remain orthogonal, and distinct half position
types remain orthogonal.  Individual color multiplicity vectors can
mix; they were kept together in \(Q\), so no orthogonality among them
is used.  Applying \eqref{s:lowering-one-step} successively to each
projected component gives
\begin{equation}\label{s:lowering-ordered-lower}
 U_{i,s}(v)\ge
 \sum_{\substack{l_0,l_1,\mu_0,\mu_1\\h_i\ge s}}
 (h_i)_s
 \prod_{t=0}^{s-1}(m-l_i-j_i+t+1)\,
 \|Q_{l_0,l_1;\mu_0,\mu_1}v\|^2 ,
\end{equation}
where \((h)_s=h(h-1)\cdots(h-s+1)\).  At the \(t\)-th stage the
position type is still \(\mu_i\), while the count is \(l_i-t\);
these are exactly the two factors from
\eqref{s:lowering-one-step}.  The sum is over ordered choices, so
there is no division by \(s!\).
Both \(l_i,j_i\le k\), and every second factor
in \eqref{s:lowering-ordered-lower} is at least \(m/2\).
Since \((h_i)_s\ge s!\) for \(h_i\ge s\), we obtain
\begin{equation}\label{s:lowering-event-lower}
 U_{i,s}(v)\ge(m/2)^s s!\Pr(h_i\ge s).
\end{equation}
Thus \(U_{i,s}(v)\ge(c_2m)^s s!\Pr(h_i\ge s)\) with \(c_2=1/2\).

\paragraph{The half-difference and pair count.}
Let
\[
 G_c=L_{X_0,c}+L_{X_1,c},\qquad
 \Delta_c=L_{X_0,c}-L_{X_1,c}
          =\sum_{p=1}^m\delta_{p,c},\qquad
 \delta_{p,c}=a_{(p,0),c}-a_{(p,1),c}.
\]
All these lowerings commute, and \(G_cv=0\).  Expanding
\(L_{X_0,c}=(G_c+\Delta_c)/2\) or
\(L_{X_1,c}=(G_c-\Delta_c)/2\), every term containing a \(G_c\)
vanishes on \(v\).  Consequently,
\begin{equation}\label{s:lowering-half-scale}
 L_{X_i,c_s}\cdots L_{X_i,c_1}v
  =(-1)^{is}2^{-s}\Delta_{c_s}\cdots\Delta_{c_1}v,
 \qquad
 U_{i,s}(v)=4^{-s}\sum_{c_1,\ldots,c_s}
             \|\Delta_{c_s}\cdots\Delta_{c_1}v\|^2 .
\end{equation}

We give the pair calculation including repeated letters.  For one pair
write \(\sigma\) for its swap and \(P_{\ge1},P_2\) for its projections
onto at least one and exactly two excitations.  On its excitation
sectors \(0,1,2\), respectively,
\[
 K_1:=\sum_c\delta_c^*\delta_c
    =0\oplus(I-\sigma)\oplus2I,\qquad K_1\preceq2P_{\ge1}.
\]
Moreover
\[
 \delta_b\delta_c
   =-\bigl(a_{(p,0),b}a_{(p,1),c}
           +a_{(p,0),c}a_{(p,1),b}\bigr).
\]
It vanishes off the two-excitation sector.  On the full pair space,
\[
 K_2:=\sum_{b,c}(\delta_b\delta_c)^*(\delta_b\delta_c)
        =2(I+\sigma)P_2\preceq4P_2.
\]
In particular \(\delta_c^2=-2a_{(p,0),c}a_{(p,1),c}\) is included
with its correct square.  Every product of three lowerings in the same
pair is zero because one of its two sites must be lowered twice.

Expand an ordered product of \(s\) of the sums \(\Delta_c\).
A nonzero term assigns the \(s\) ordered positions to a set \(A\) of
pairs used once and a disjoint set \(Z\) of pairs used twice.
Put \(a=|A|\), \(z=|Z|\), so \(a+2z=s\).  For fixed \(A,Z\) there
are exactly
\[
 N_{a,z}=\frac{s!}{2^z}
\]
such assignments; the ordered positions remain distinct even when
their letters agree.  If \(T_{f,\mathbf c}\) denotes the product for
one assignment \(f\), then
\[
 \sigma_p T_{f,\mathbf c}v
     =(-1)^{\mathbf1_{\{p\in A\}}}T_{f,\mathbf c}v .
\]
Indeed \(\sigma_p\delta_{p,c}\sigma_p=-\delta_{p,c}\),
it commutes with the other pair lowerings, and \(\sigma_pv=v\).
The pair swaps are commuting self-adjoint unitaries.  Thus terms with
different singleton sets \(A\) lie in orthogonal joint eigenspaces.
For fixed \(s\), the size of \(A\) also fixes \(z\).

Let
\[
 P_{A,Z}=\prod_{p\in A}P_{p,\ge1}\prod_{p\in Z}P_{p,2}.
\]
Summing the ordered letters for one assignment, the disjoint tensor
factors and the two local squares above give
\[
 \sum_{\mathbf c}\|T_{f,\mathbf c}v\|^2
 =\left\langle v,
        \prod_{p\in A}K_{1,p}\prod_{p\in Z}K_{2,p}v\right\rangle
 \le2^a4^z\|P_{A,Z}v\|^2=2^s\|P_{A,Z}v\|^2.
\]
For a fixed \(A\) there are \(\binom{m-a}{z}\) choices of \(Z\).
Cauchy--Schwarz over those choices and over their \(N_{a,z}\)
assignments costs \(\binom{m-a}{z}N_{a,z}\); after the letter sum,
the sum over assignments contributes the second \(N_{a,z}\).
Using the orthogonality between different \(A\)'s and then
\eqref{s:lowering-half-scale}, we conclude that
\begin{equation}\label{s:lowering-pair-upper}
 U_{i,s}(v)\le
 2^{-s}(s!)^2
 \sum_{a+2z=s}4^{-z}\binom{m-a}{z}
       \sum_{\substack{|A|=a,\ |Z|=z\\A\cap Z=\varnothing}}
                  \|P_{A,Z}v\|^2 .
\end{equation}
In this display the inner sum is multiplied by the binomial factor only
once; it charges interference between the possible double sets \(Z\)
with the same parity set \(A\).

For a word \(w\), let \(t(w)\) be its number of doubly excited pairs
and \(b(w)\) its number of pairs with at least one excitation.
The total-\(k\) condition gives \(b(w)+t(w)=k\).  All \(P_{A,Z}\)
are diagonal word projections, and the number that retain \(w\) is
\(\binom{t(w)}z\binom{b(w)-z}a\): choose \(Z\) among the doubles,
then \(A\) among the other occupied pairs.  Therefore
\[
 \sum_{\substack{|A|=a,\ |Z|=z\\A\cap Z=\varnothing}}
       \|P_{A,Z}v\|^2
 =\mathbb E_{w\sim|\langle w,v\rangle|^2}
       \left[\binom{t(w)}z\binom{b(w)-z}a\right]
 \le\frac{k^{a+z}}{a!\,z!}.
\]
We use the value zero for infeasible binomial coefficients.
The squared half factor \(4^{-s}\) has already entered
\eqref{s:lowering-pair-upper}. Inserting the preceding occupancy bound
there and using \(\binom{m-a}{z}\le\binom mz\) and
\(2^{-s}4^{-z}\le1\) gives the coarser ordered-letter bound with unit
coefficient \(C_2=1\):
\[
 U_{i,s}(v)\le (s!)^2
       \sum_{a+2z=s}\binom mz\frac{k^{a+z}}{a!\,z!}.
\]
To obtain the tail estimate, return to the more precise factors in
\eqref{s:lowering-pair-upper}. Set \(x=k/m\le1/4\).
Dividing that bound by \eqref{s:lowering-event-lower} and using
\(\binom{m-a}{z}\le m^z/z!\) gives
\[
 \Pr(h_i\ge s)
 \le s!\sum_{a+2z=s}\frac{x^{a+z}}{4^z a!(z!)^2}
 \le\sum_{a+2z=s}\frac{s!}{a!(2z)!}x^a(\sqrt{x})^{2z}
 \le(x+\sqrt{x})^s
 \le(4x)^{s/2}.
\]
The middle estimate uses \(\binom{2z}{z}\le4^z\), and the next sum
contains only the even terms of the binomial expansion.  For \(s>k\)
the event is empty, so the conclusion holds for every integer \(s\ge1\).
Finally put \(u=h_0+h_1\).  If \(u\ge r\), for any real \(r\ge1\),
then some \(h_i\ge\lceil r/2\rceil\).  Since \(4x\le1\),
\begin{equation}\label{ten:lowering:eq:2}
 \Pr(u\ge r)\le2(4k/m)^{r/4}
              \le C_1(C_1k/m)^{r/4},
 \qquad C_1=4,\quad r\ge1 .
\end{equation}
All constants are independent of the alphabet, the shape and the
multiplicity spaces.
On the remaining parent range \(k/m>1/4\), the final bound in
\eqref{ten:lowering:eq:2} is automatic, since
\(C_1(C_1k/m)^{r/4}=4(4k/m)^{r/4}\ge4\).
Thus that bound holds throughout \(2\le k\le n/2\).

For completeness, defect \(1\) still vanishes exactly.  On the
one-excitation space each layer averages the indicated bit of the
position, so their product is the projection onto constant position
functions.  Its orthogonal complement is the position type
\((n-1,1)\), and hence \(B_n\) is zero on that type.

\subsection{Removing prefactor losses}
The tail estimate now closes the induction without accumulating a loss
at each split.  Write
\[
 W_n(\lambda)=\|B_n\|_\lambda^2,
 \qquad h_*=\frac1{20},\qquad c_{\rm s}=\frac1{64}.
\]
We will choose an absolute $M_*\ge1$ and prove the stronger bound
\begin{equation}\label{s:lowering-margin-invariant}
 W_n(\lambda)\le h_*\left(\frac{M_*k}{n}\right)^{c_{\rm s}k}
 \qquad(2\le k\le n/2).
\end{equation}
This is the potential
$V(n,k)=k\log(n/k)-Ak+B\mathbf1_{\{k>0\}}$ written with
$M_*=e^A$ and $h_*=e^{-c_{\rm s}B}$; as usual $V(n,0)=0$.
The fixed factor $h_*$ supplies a margin when both halves retain
positive defect.

If $W_n(\lambda)=0$, the estimate is immediate; otherwise use the
component law above.  Suppose the proposed right side in
\eqref{s:lowering-margin-invariant} is less than one; otherwise
contractivity proves the estimate.  Put $y=k/n$ and
\[
 D=\frac{(l_0-k/2)^2}{k},\qquad
 u=k-j_0-j_1,\qquad b=\#\{i:j_i>0\}.
\]
We will choose $M_*$ so that this nontrivial case has $y<\delta_*$,
where $0<\delta_*\le1/8$.  Then $k<m/4$, so every child defect
$j_i\ge2$ lies in the range of the induction hypothesis.
The trivial child has norm one, and a child with defect one has norm
zero.  For a component with neither child defect equal to one, the
product of the two child bounds divided by the proposed parent bound
is exactly
\begin{equation}\label{s:lowering-induction-ratio}
 h_*^{b-1}(M_*y)^{-c_{\rm s}u}
 \exp\left(c_{\rm s}\sum_{i=0}^1j_i\log(2j_i/k)\right),
\end{equation}
with zero summands at $j_i=0$.

The exponent in this ratio has a simple uniform bound.  Using
$\log z\le z-1$ and $j_0+j_1=k-u$ gives
\begin{align*}
 \sum_i j_i\log(2j_i/k)
 &\le \frac{2(j_0^2+j_1^2)}k-(k-u)\\
 &=\frac{(j_0-j_1)^2}{k}-\frac{u(k-u)}k
 \le\frac{(j_0-j_1)^2}{k}.
\end{align*}
Since $j_0-j_1=(l_0-l_1)-(h_0-h_1)$ and
$|h_0-h_1|\le u\le k$, the last expression is at most
\begin{equation}\label{s:lowering-induction-cost}
 8D+2u.
\end{equation}
The split tail in \eqref{s:lowering-split-law} implies
$\Pr(D>z)\le2e^{-2z}$ for $z\ge0$.  Consequently, for $0\le t<2$,
\[
 \mathbb E e^{tD}
 =1+t\int_0^\infty e^{tz}\Pr(D>z)\,dz
 \le1+\frac{2t}{2-t}.
\]
In particular $\mathbb E e^{16c_{\rm s}D}\le9/7<2$.

On the event $u=0$, the child defects equal their excitation counts.
If $b=1$, all excitations lie in one half, so
\eqref{s:lowering-split-law} and the exact exponent
$\sum_i j_i\log(2j_i/k)=k\log2$ bound this part of the expectation
of \eqref{s:lowering-induction-ratio} by
\[
 2^{1-k}2^{c_{\rm s}k}
 \le2^{2c_{\rm s}-1}<\frac35.
\]
If $b=2$, \eqref{s:lowering-induction-cost} bounds its contribution by
$h_*\mathbb E e^{8c_{\rm s}D}\le2h_*=1/10$.

It remains to bound positive loss.  Since $M_*\ge1$ and
$h_*^{b-1}\le h_*^{-1}$, Cauchy--Schwarz and
\eqref{ten:lowering:eq:2}, in the form
$\Pr(u\ge r)\le4(8y)^{r/4}$, bound the total contribution from
$u\ge1$ by
\begin{align*}
 \frac1{h_*}\sum_{r\ge1}y^{-c_{\rm s}r}e^{2c_{\rm s}r}
       \mathbb E[\mathbf1_{\{u=r\}}e^{8c_{\rm s}D}]
 &\le\frac{2\sqrt2}{h_*}\sum_{r\ge1}
   \left(e^{2c_{\rm s}}8^{1/8}y^{1/8-c_{\rm s}}\right)^r.
\end{align*}
Here $1/8-c_{\rm s}=7/64>0$.  Choose $0<\delta_*\le1/8$ so that
the last geometric series, with $y$ replaced by $\delta_*$, is at
most $1/5$.  Then choose
\[
 M_*\ge\delta_*^{-1}h_*^{-1/(2c_{\rm s})}.
\]
If $y\ge\delta_*$ and $k\ge2$, this choice makes
$h_*(M_*y)^{c_{\rm s}k}\ge1$.  Thus every nontrivial case indeed
has $y<\delta_*$, as required for the tail estimate and child range.

The three contributions to the ratio are now less than
$3/5+1/10+1/5=9/10$.  Components with a defect-one child contribute
zero.  Applying \eqref{s:lowering-backward-recursion} proves the
inductive step.  The range $2\le k\le n/2$ is empty at $n=2$,
and all cases with a proposed bound at least one were already covered
by contractivity.  Strong induction on $d$ proves
\eqref{s:lowering-margin-invariant}.  Dropping $h_*<1$ gives
\eqref{ten:lowering:eq:1} with the absolute choices
$c=c_{\rm s}>0$ and $C=c_{\rm s}\log M_*$.

\subsection{Signed types and the inherited hook budget}

We now prove Proposition~\ref{ten:lowering:main}. The norm estimate
settles small first-row defects. At the other shapes we will split a
sweep into a rectangle, weight each local signed type by its entropy,
and interpolate a trace bound with an operator bound. The following
comparisons connect those local types with the subdiagrams in
$L_n(\lambda)$.

First, for a $q$-hook diagram $\eta$ of size at most $n$,
\begin{equation}\label{s:lowering-frobenius-dimension}
 \log D_\eta=F(\eta)+O(q^2\log(n+2)).
\end{equation}
Here is the hook-formula comparison, including the choice of Frobenius
parts. Put $A_i=a_i-1$, $B_i=b_i$, with $r\le q$ the Durfee index.
The Frobenius form of the hook formula is
\[
 D_\eta=|\eta|!\,
 \frac{\prod_{i<j}(A_i-A_j)(B_i-B_j)}
      {\prod_i A_i!B_i!\prod_{i,j}(A_i+B_j+1)}.
\]
It follows by splitting the hook product at the Durfee square;
compare the row form in~\cite[\S4.17, formula (5) and Theorem 4.53]{etingof}.
All the displayed cross and Vandermonde factors are positive integers
at most $2n+1$, and there are $O(q^2)$ of them. Replacing $A_i!$
by $a_i!$ costs another $O(q\log(n+2))$ in logarithms. Finally the
multinomial logarithm for the parts $a_i,b_i$ differs from their
entropy by $O(q\log(n+2))$. This proves
\eqref{s:lowering-frobenius-dimension}, with the empty shape treated
as dimension one and entropy zero.

Use now the signed alphabet with $q$ even and $q$ odd colors from
Section~\ref{sec:static}. A compact type is a pair
$\gamma=(\gamma^+,\gamma^-)$; write $F_\gamma$ for the entropy of
its concatenated parts. Lemma~\ref{lem:signed-hook-carriers} gives
all the needed conversion facts. The occurring permutation types are
exactly the $q$-hook shapes, their full signed-tensor multiplicities
are $\exp(O(q^2\log(n+2)))$, and, whenever $\eta$ occurs with
compact type $\gamma$,
\[
 F_\gamma-O(q^2\log(n+2))\le\log D_\eta\le F_\gamma.
\]
Combining this with \eqref{s:lowering-frobenius-dimension} yields
\[
 F_\gamma=F(\eta)+O(q^2\log(n+2)).
\]
Thus both the compact carrier and the permutation multiplicities
fit the same error, and every permitted hook can be realized by
at least one compact type.


The parent uses $q_n$ colors, whereas a child moment at size $m$
is stated with the smaller hook parameter $q_m$. The next comparison
justifies using the larger inherited types in that child estimate.

\begin{lemma}[Passing to an inherited hook]\label{s:lowering-inherited-hook}
Suppose, at a sufficiently large dyadic size $m$, positive numbers
$E_\mu$ satisfy
\begin{equation}
 \log E_\mu\le \beta F(\eta_0)+G_m(m-|\eta_0|)
 \quad\text{for every $q_m$-hook $\eta_0\subseteq\mu$.}
 \label{ten:lowering:eq:4}
\end{equation}
Then every subdiagram $\eta\subseteq\mu$, of size $m-l$, satisfies
\begin{equation}
                 \log E_\mu\le\beta F(\eta)+G_m(l).
                                                                  \label{ten:lowering:eq:5}
\end{equation}
The same conclusion is automatic when $E_\mu=0$, with logarithm
$-\infty$.
\end{lemma}
\begin{proof}
Intersect $\eta$ with the $q_m$-hook, obtaining $\eta_0$ by deleting
$x$ boxes. This removes precisely its Frobenius parts of index
greater than $q_m$, each of size at most $|\eta|/q_m$.
Separating the removed parts from the retained ones gives
\[
 F(\eta)-F(\eta_0)\ge x(\log q_m-1).
\]
On each interval of differentiability, $G_m'(t)\le\alpha\log m+1$;
the function is continuous at its clipping threshold. Hence
\[
 \beta F(\eta_0)+G_m(l+x)
 \le \beta F(\eta)+G_m(l)
       +x\bigl[\alpha\log m+1-\beta(\log q_m-1)\bigr].
\]
The bracket is nonpositive for all sufficiently large $m$, since
$\beta\delta>\alpha$. Apply the hypothesis to $\eta_0$.
\end{proof}

\subsection{The range requiring a moment induction}

The moment bound \eqref{ten:lowering:eq:6} holds automatically on
the trivial module. At any fixed finite collection of sizes,
Lemma~\ref{lem:finitegap} allows one exponent to make every
nontrivial weighted moment at most one. We will choose that
exponent after fixing the threshold for the large-size argument.

Let us also dispense with some representations. Diagrams of height \(>n/2\) don't survive even one layer (invariants under disjoint switches require dominance over the content with \(n/2\) twos, by the Young/Pieri rule). For \(k\le n^{1-\alpha/16}\) (first row defect), \eqref{ten:lowering:eq:6} at large sizes follows already with an absolute power taken sufficiently large, by \eqref{ten:lowering:eq:1}, since \(D_\lambda\le n^{2k}\). All constants, even ones depending on our small fixed parameters, are meant to be absolute. For remaining shapes write \(S=\log D_\lambda\); we note
\[
                            S\gtrsim n^{1-\alpha/16}.
\]
Here is one verification. If height and width are both \(< n/10\), use the hook bound. Otherwise by transposing take a dimension at least \(n/10\) as width. Below the first row still lie at least \(c n^{1-\alpha/16}\) cells (using the height restriction in the transposed case). Take a subdiagram consisting of the first row shortened if necessary but still length \(\lfloor n/10\rfloor\), and \(l=\lfloor c' n^{1-\alpha/16}\rfloor\) cells below, \(c'\) small. It has exponentially many tableaux on scale \(l\): for instance fill the first \(l\) places of row one, then interleave the rest of that row freely with a tableau below. Branching proves the bound.

Moreover on these remaining shapes
\begin{equation}
                              L_n(\lambda)<\tfrac12 S                       \label{ten:lowering:eq:7}
\end{equation}
at large sizes. Take the intersection with the \(q_n\)-hook for \(\eta\). Its dimension is at most \(D_\lambda\). If \(t=n-|\eta|\) is below \(n^{1-\alpha/4}\), \eqref{ten:lowering:eq:7} is immediate. Otherwise the \(t\) cells cut away themselves give a diagram (after translation) contained in \(\lambda\), of width and height at most \(n/q_n\), giving \(S\ge (\delta/2)t\log n\) by the hook bound. This proves \eqref{ten:lowering:eq:7}.

\subsection{Weighted interpolation on a rectangle}
For the classical Schatten three-lines interpolation principle used here,
see the Stein--Hirschman formulation in
\cite[Section 3.1, Theorem 3.1]{SutterBertaTomamichel2017}; uniform boundary
bounds give the form needed below. The weighted family and its estimates
are developed here.

Split the chronological coordinate order into two nearly equal groups.
The first gives rows of size $m_1$ and the second columns of size $m_2$,
where $m_1m_2=n$ and both sizes are within a factor two of $\sqrt n$.
Write
\[
 B_n=CR,
\]
where $R$ is the product of the independent row sweeps and $C$ the
product of the independent column sweeps. Take
$p=\max(p_{m_1},p_{m_2})$; the child bounds remain valid at this
larger power.

Choose a minimizing hook $\eta$ in \eqref{ten:lowering:eq:6} and put
$t=n-|\eta|$. Let $V=E\oplus O$ have $q=q_n$ even and $q$ odd
colors, and let $H$ be the even space with basis $h_1,\ldots,h_t$.
The compact group $G_q$ acts on $V$ and trivially on $H$. Let
$W\subseteq(V\oplus H)^{\otimes n}$ be the span of words in which
each marker $h_a$ appears exactly once. Choose a global compact type
$\Gamma$ that occurs with permutation shape $\eta$ on the $n-t$
unmarked sites, and let $P=P_\Gamma$. In $PW$, the target shape
$\lambda$ occurs with multiplicity at least $f^{\lambda/\eta}$,
by branching with the distinct markers. In particular $P$ reduces
both $R$ and $C$.

The positive operators needed at the middle of this product are
\[
 A=(RR^*)^{1/2},\qquad B=(C^*C)^{1/2}.
\]
To see their orientation explicitly, extend polar partial isometries
to unitaries within each irreducible block of the corresponding
subgroup algebra. Then $R=AU_R$, $C=U_CB$, and
$CR=U_CBAU_R$. These unitaries commute with $P$, so the sweep restricted to $PW$
has Schatten $p$ norm $\|BPA\|_p$. This remains true when either
sweep has a kernel.

For a local block $X$ of size $m$, let $Q_{X,\gamma}$ be its full
$G_q$-isotypic projection on $(V\oplus H)^{\otimes X}$, where
$\gamma=(\gamma^+,\gamma^-)$ has total size at most $m$ and each
partition has at most $q$ parts. The total size of $\gamma$ is the
number of clean letters. This projection includes every compatible
marker word and every placement of its letters. Define
\[
 F_X=\sum_\gamma F_\gamma Q_{X,\gamma},\qquad
 F_A=\sum_{X\text{ row}}F_X,\qquad
 F_B=\sum_{X\text{ column}}F_X,
\]
and restrict the resulting operators to $W$. Each local sweep
preserves marker contents and commutes with the color action. Thus
$F_A,F_B$ commute with their own positive operators $A,B$, respectively,
and both commute with the global compact-type projection $P$.
We interpolate a trace boundary carrying weights
$\sigma^{p/2}e^{(1-\beta)F_\gamma/2}$ with an operator boundary
carrying weights $e^{-(1-\beta)F_\gamma/(p-2)}$. Here $\sigma$
is the local singular value operator, with the orientation just
specified.

For $p>2$ put
\[
 h(z)=(1-\beta)\left(\frac z2-\frac{1-z}{p-2}\right),\qquad
 \mathcal Z(z)=B^{pz/2}e^{h(z)F_B}P e^{h(z)F_A}A^{pz/2}.
\]
At $\vartheta=2/p$ one has $h(\vartheta)=0$ and $\mathcal Z(\vartheta)=BPA$.
On both boundaries the imaginary powers lie on the outside and have
norm at most one. Zero singular spaces stay zero throughout.
Schatten interpolation from the operator norm to the Hilbert--Schmidt
norm gives
\begin{equation}\label{s:lowering-weighted-strip}
 \|BPA\|_p^p\le\|\mathcal Z(1)\|_2^2\|\mathcal Z(0)\|_\infty^{p-2}.
\end{equation}
We will choose $p$ large enough that $2(1-\beta)/(p-2)\le1$.

\subsection{Positive compression at marked sites}
We first bound the Hilbert--Schmidt norm squared. We may discard \(P_\Gamma\), by global equivariance, and are bounding the trace of the product of the two squared-weight operators. Terms in this trace must have identical hole positions between the two matrices: each hole must stay in its row and in its column. At a fixed ordered marker placement, $Q_{X,\gamma}$ restricts to the clean compact-type projection, with the same entropy weight $F_\gamma$. The following local estimate is what we need. If a block holds \(l\) holes, the trace over remaining clean letters, on diagonal compression fixing the positions of these holes, of the positive squared-weight operator in the block (all types) is at most
\begin{equation}
         (m)_l^{-1}\exp\{G_m(l)+O(q^2\log(n+2)+\sqrt m)\},                 \label{ten:lowering:eq:8}
\end{equation}
where falling factorials are used.

To see this, decompose the local singular power by permutation shapes \(\mu\) at size \(m\). For a positive matrix coefficient operator supported on \(\mu\), with matrix there \(T\ge0\), the diagonal compression just uses its convolution coefficients restricted to the subgroup on \(m-l\) positions. Signs from the tensor convention are harmless or are absorbed by putting the hole positions last (signed tensor reordering). In Fourier terms, on a shape \(\eta'\) of this subgroup the result of restriction of coefficients is positive, of trace at most
\[
               \frac{D_\mu\operatorname{Tr}T}{(m)_l D_{\eta'}},
\]
and requires \(\eta'\subset\mu\). Indeed restriction with normalized convolution scaling takes the partial trace on the multiplicity factor of the \(\eta'\) part of \(T\), with exactly the displayed dimension and factorial factor, by Schur orthogonality. Apply this in the clean module, using \(T=\sigma^p|_\mu\), and take trace on a given clean type, multiplying by its squared color weight factor at this boundary. Now \eqref{ten:lowering:eq:5}, the induction, and the entropy and dimension approximations \(F(\eta')=F_\gamma+O(q^2\log(n+2)),\ \log D_{\eta'}=F_\gamma+O(q^2\log(n+2))\), pay the entire color weight. Summing types and \(\mu\)'s with their possible clean multiplicities gives \eqref{ten:lowering:eq:8}; we have used the standard \(\exp(O(\sqrt m))\) partition-number bound. For disjoint local blocks we can multiply the bounds. More explicitly in signed tensor space, one reorders sites to group blocks to see the tensor product, and on fixed hole positions the reordered positive operators have the same trace. For the row diagonal compression and column diagonal compression, bound trace of their product by product of their traces.

Write
\[
             E=n^{3/4} \log^2(n+2)+n^{1/2+2\delta}\log^2(n+2).
\]
A uniform bound for the log of this boundary Hilbert--Schmidt norm squared is
\begin{equation}
 t(\alpha\log n-2\log(n/t))+O\big(t+E+n^{1-\alpha/3}\log(n+2)\big),        \label{ten:lowering:eq:9}
\end{equation}
with the main term zero for \(t=0\). We detail summation. Replace falling factorial reciprocals by \(m^{-l}\) costing \(O(t)\) in the log per side. The number of placed distinct holes is at most \(n^t\), canceling the \(m\)-powers, so average the remaining exponential over \(t\) uniform distinct sites. In \(G_m(l)\) we can remove the truncation: total error is \(O(n^{1-\alpha/3}\log(n+2))\), using \(m\asymp\sqrt n\) and that discrepancies occur only at \(l\le m^{1-\alpha}\). Let \(x,y\) be empirical row and column proportions of the hole sites. The sum of the log terms is then the main term in \eqref{ten:lowering:eq:9}, plus \(t\) times the two relative entropies of \(x,y\) versus uniform. The latter exponential factors cost only \(O(t+E)\) in the log. Indeed replacing distinct sampling by independent samples costs at most \(e^{O(t)}\), and the two marginal histograms are then independent with probability for given histograms bounded by the inverse entropy factors; there are at most \((t+1)^{m_1+m_2}\) choices. This proves \eqref{ten:lowering:eq:9}, the \(E\) term also counting all the local errors.

\subsection{Density normalization at the operator boundary}
For the operator boundary we explain a whitening estimate. It will give, in logs \emph{after raising to \(p-2\)}, at most
\begin{equation}
                  -(1-\beta) F_\Gamma + O(t+E).                            \label{ten:lowering:eq:10}
\end{equation}
It suffices to work at magnitude weights \(e^{-F_\gamma/2}\), then steepness can be reduced to \((1-\beta)/(p-2)\le1/2\), by operator-norm interpolation between these weights and identity bounds, scaling the log estimate. Outside supported projections may only decrease the bound. Aim then for log norm \(-F_\Gamma/2+O(t+E)\).

We first construct one positive mixture for the entire local weight
$e^{-F_X}$. Fix a clean type $\gamma$ and write
\[
 s=|\gamma^+|+|\gamma^-|,\qquad
 \mathcal U_\gamma=\mathcal U_{\gamma^+}\otimes\mathcal U_{\gamma^-},
 \qquad u_\gamma=\dim\mathcal U_\gamma.
\]
For $s>0$, take the even density $D_\gamma$ on $V$ with eigenvalues
$\gamma_i^+/s$ on $E$ and $\gamma_j^-/s$ on $O$, padded by zeroes.
For $s=0$ choose any even density. For $g\in G_q$ put
\[
 \widehat D_{\gamma,g}=(gD_\gamma g^*)\oplus I_H.
\]
Only the clean restriction has trace one; each marker has eigenvalue
one.

Here is why this single choice controls every marker fiber with type
$\gamma$. Fix the clean positions and the letters at all other
positions. Reordering the clean factors gives the decomposition
$\mathcal M_\gamma\otimes\mathcal U_\gamma$ of
\eqref{eq:signed-sector-decomposition}. On this space
$(D_\gamma\oplus I_H)^{\otimes X}$ acts as
$I_{\mathcal M_\gamma}\otimes\pi_\gamma(D_\gamma)$, where
$\pi_\gamma$ is the polynomial color action. Summing these orthogonal
spaces over clean positions and marker words only enlarges the
multiplicity factor. Denote this full multiplicity space by
$\mathcal N_{X,\gamma}$. On the whole $Q_{X,\gamma}$ range the action is
\[
 I_{\mathcal N_{X,\gamma}}\otimes\pi_\gamma(D_\gamma).
\]
Thus Haar averaging is the same scalar on every copy of
$\mathcal U_\gamma$, independently of marker labels and placements.
Lemma~\ref{lem:signed-type-density} identifies that scalar as at least
$e^{-F_\gamma}/u_\gamma$. The average preserves every compact type
and is positive there, so on the full local tensor space
\[
 e^{-F_\gamma}Q_{X,\gamma}\preceq
 u_\gamma\int_{G_q}\widehat D_{\gamma,g}^{\otimes X}\,dg.
\]
The same assertion includes the empty clean type: its carrier is
one-dimensional and all its factors are markers.

We can now add these inequalities over clean types. There are at most
$(m+1)^{2q}$ pairs of partitions of total size at most $m$, and
$u_\gamma\le(m+1)^{q(q-1)}$. Consequently
\[
 e^{-F_X}\preceq
 \sum_\gamma u_\gamma\int_{G_q}\widehat D_{\gamma,g}^{\otimes X}\,dg,
 \qquad
 c_X:=\sum_\gamma u_\gamma\le(m+1)^{q^2+q}.
\]
This sum charges each clean type once, with all its marker fibers
already included in $Q_{X,\gamma}$. Tensor these positive inequalities
over the rows or the columns. Every factor preserves the sector $W$,
so restriction to $W$ gives positive mixtures dominating $e^{-F_A}$
and $e^{-F_B}$. After dividing by their masses, the index measures
are probabilities, and the two masses satisfy
\[
 \log c_A+\log c_B
 \le (q^2+q)(m_1+m_2)\log(n+2).
\]
Apply Lemma~\ref{ten:positive-mixtures} to these two entropy weights,
with middle operator $P_\Gamma$. It bounds
$\|e^{-F_B/2}P_\Gamma e^{-F_A/2}\|$ by
$\sqrt{c_Ac_B}$ times the supremum of
\[
 \|T_{\rm col}^{1/2}P_\Gamma T_{\rm row}^{1/2}\|,
\]
where $T_{\rm col}$ and $T_{\rm row}$ are sitewise tensor products
of extended densities, one clean density for each column or row,
repeated along that line. This comparison acts on arbitrary vectors
in $W$. Its logarithmic mass cost is
$O(q^2(m_1+m_2)\log(n+2))$.

For one pair of mixture terms, write $R_i$ for its clean row
densities and $C_j$ for its clean column densities. Put
$Z=m_2^{-1}\sum_{i=1}^{m_2}R_i$, a positive trace-one matrix on $V$.
Use inverse powers on its support, and extend all powers by identity
on $H$. In the middle write
\[
       P_\Gamma=(P_\Gamma (Z^{1/2})^{\otimes n})
                      (Z^{-1/2})^{\otimes n}
\]
on the support needed when applied to the row product. The same
identity-on-multiplicity decomposition applies to the global type
$\Gamma$, with $n-t$ clean letters. Its first factor therefore has
norm at most $e^{-F_\Gamma/2}$ by
Lemma~\ref{lem:signed-type-density}. Apply
Lemma~\ref{s:compact-coefficient-extraction} to the carrier matrix
$A=\pi_\Gamma(Z^{1/2})$. The factor
$P_\Gamma(Z^{1/2})^{\otimes n}$, zero on all other compact types,
is an integral of global color actions with absolute coefficient
integral at most
\[
 (\dim\pi_\Gamma)e^{-F_\Gamma/2}
 \le e^{-F_\Gamma/2}\exp(O(q^2\log(n+2))).
\]
The extended densities, the global color actions and $P_\Gamma$
preserve each ordered marker placement. Their product is therefore
an orthogonal direct sum over placements, and its operator norm is
the supremum of the norms on those summands. On clean sites we need
the tensor over occupied cells (non-hole sites) of
\[
              C_j^{1/2}\, U Z^{-1/2}\, R_i^{1/2}
\]
with $U\in G_q$. The sum over all $n$ cells of their squared
Frobenius norms is at most $n$: summing the whitened row factors
gives $m_2$ times the support projection of $Z$, and each $C_j$ has
trace one. Lemma~\ref{ten:missing-cells} therefore bounds the product
of the one-cell norms on any $n-t$ occupied cells by $e^{t/2}$.
The empty product is one when all sites are holes. This bound is
uniform over marker placements, so taking their supremum adds no
factor. Together with the mixture masses and the global carrier
factor, it proves
\[
 \log\|e^{-F_B/2}P_\Gamma e^{-F_A/2}\|
 \le-F_\Gamma/2+O(t+E).
\]
To obtain the precise operator boundary, interpolate
$e^{-zF_B/2}P_\Gamma e^{-zF_A/2}$ between this estimate and the
norm-one boundary at $\Re z=0$. At
$\vartheta=2(1-\beta)/(p-2)\le1$ the result is
\[
 \log\|\mathcal Z(0)\|\le
 \frac{-(1-\beta)F_\Gamma+O(t+E)}{p-2}.
\]
Raising to the power $p-2$ in \eqref{s:lowering-weighted-strip}
therefore gives \eqref{ten:lowering:eq:10} with an error constant
independent of $p$. This order of interpolation permits the finite-base
power to be chosen after all size thresholds.

\subsection{Closing the moment induction}
Putting together \eqref{ten:lowering:eq:9}, \eqref{ten:lowering:eq:10} and interpolation, and using multiplicity at least \(f^{\lambda/\eta}\), we obtain
\begin{equation}
\begin{split}
 \log( D_\lambda\operatorname{Tr}_\lambda |B_n|^p )
 \le {}&\log(D_\lambda/f^{\lambda/\eta})-(1-\beta) F_\Gamma\\
 &+t(\alpha\log n-2\log(n/t))
                   +O(t+E+n^{1-\alpha/3}\log(n+2)).
\end{split}                                                               \label{ten:lowering:eq:11}
\end{equation}
We address a dimension detail here: \(D_\lambda\) compared to \(D_\eta f^{\lambda/\eta}\) need not be estimated by counting orderings through \(\eta\) alone without any loss. We need an upper bound at cost \(\binom{n}{t}\) times at most exponential in \(t\). In fact
\[
                       D_\lambda\le \binom{n}{t} D_\eta f^{\lambda/\eta}.
\]
A standard tableau determines the subset of its entries in \(\eta\), their ordering there and the ordering on the skew complement, proving exactly this bound. Hence \eqref{ten:lowering:eq:11} gives
\begin{equation}
               \log( D_\lambda\operatorname{Tr}_\lambda |B_n|^p )
                  \le L_n(\lambda)+O(t+E+n^{1-\alpha/3}\log(n+2)).
                                                                         \label{ten:lowering:eq:12}
\end{equation}
It was relevant both to transfer the clean dimension through the middle and to have the negative log-sparsity term in \eqref{ten:lowering:eq:3}, so that two contributions at subsizes can accommodate this binomial cost.

Write the error in \eqref{ten:lowering:eq:12} as $\varepsilon_n$.
For the remaining shapes,
\[
 \varepsilon_n\le K\frac{S}{\log n},\qquad S=\log D_\lambda,
\]
with $K$ independent of $p$. To see this for its $O(t)$ term,
if $t<n^{1-\alpha/8}$ use $S\gtrsim n^{1-\alpha/16}$.
Otherwise the minimizing cost contains
\[
 G_n(t)\ge\frac{7\alpha}{8}t\log n,
 \qquad G_n(t)\le L_n(\lambda)<S/2,
\]
so $t\le4S/(7\alpha\log n)$. The terms
$E+n^{1-\alpha/3}\log(n+2)$ are $o(S/\log n)$ in both cases.

Take a size threshold above which $K/\log n\le1/4$. Equation
\eqref{ten:lowering:eq:12} and $L_n(\lambda)<S/2$ imply, for each
singular value $s$ of $B_n|_{V_\lambda}$,
\[
 s^p\le\operatorname{Tr}_\lambda|B_n|^p
 \le\exp\{L_n(\lambda)+\varepsilon_n-S\}
 \le e^{-S/4}.
\]
Consequently, for $A\ge4K$,
\[
 \log\!\left(D_\lambda\operatorname{Tr}_\lambda
       |B_n|^{p(1+A/\log n)}\right)
 \le L_n(\lambda)+\varepsilon_n-\frac{AS}{4\log n}
 \le L_n(\lambda).
\]
This is the exponent increase needed for the parent estimate.

Choose an integer $d_0\ge2$ large enough for all preceding
large-size and child-size comparisons. The constants $A,d_0$ depend
only on the fixed budget parameters; the scaled operator-boundary
estimate \eqref{ten:lowering:eq:10} keeps them independent of the
base exponent. Now choose one $P_*\ge8$ that meets the sparse
requirement and makes every nontrivial weighted moment at sizes
$2^d$, $1\le d\le d_0$, at most one, using
Lemma~\ref{lem:finitegap}. The trivial moment is one and every
candidate in $L_n$ is nonnegative. Define
\[
 \begin{aligned}
 p_{2^d}&=P_* &&(1\le d\le d_0),\\
 p_{2^d}&=\left(1+\frac{A}{d\log2}\right)
       \max\{p_{2^{\lfloor d/2\rfloor}},
              p_{2^{\lceil d/2\rceil}}\} &&(d>d_0).
 \end{aligned}
\]
The preceding estimates prove the induction with these powers.
Along a branch of rounded halvings, the sum of reciprocal bit
lengths above $d_0$ is at most $2/d_0$. Therefore
\[
 p_{2^d}\le P_*\exp\!\left(\frac{2A}{d_0\log2}\right)
 \qquad(d\ge1).
\]
This proves the uniform bound on the exponents and completes
Proposition~\ref{ten:lowering:main}.

\subsection{Amplification to full-deck mixing}
Finally, the estimates imply TV convergence in constantly many butterflies. To state the details, at sparse levels \eqref{ten:lowering:eq:1} in the range used before and the elementary dimension bound show, on increasing a fixed absolute power if necessary, that the sum of \(D_\lambda\) times the power traces off the trivial diagram in this range tends to zero. Here one may sum at defect \(k\) over at most the partition number of \(k\) tails, and get for example \(n^{-2k}\) bounds per shape; defect 1 vanished exactly. For remaining levels \eqref{ten:lowering:eq:6}, \eqref{ten:lowering:eq:7} give the same conclusion with a sufficiently large fixed power, by the lower bound on log dimensions, height cutoff and partition count. Take the power even, \(2a\); increasing powers throughout is harmless. The density relative to uniform after \(a\) complete butterflies has squared \(L^2\)-distance at most this vanishing sum, by Plancherel (products of \(a\) butterfly matrices in Fourier space have Hilbert--Schmidt norm bounded by the product of Schatten-\(2a\) norms). This holds from any starting order, hence gives the required upper bound.
